\part{The global BFJ theory}
\label{part:bfj}

This part develops the global pluripotential theory initiated by
Boucksom--Favre--Jonsson over a nontrivially valued field.  It begins under the
original geometric hypotheses: a smooth projective variety over a complete
discretely valued field of residue characteristic zero, together with an
ample line bundle.  Compactness, decreasing regularization, envelopes,
Monge--Amp\`ere measures, capacity, finite energy, and the variational equation
are developed in that setting, with the algebraizability, measure-support
and continuity hypotheses recorded at each theorem. Later chapters separate hypotheses removed
by subsequent work from those that remain essential or open.

\paragraph{Prerequisites and support.}
\label{front:part3-prerequisites-and-support}
The indispensable preparation is the model-skeletal chain
\Cref{chap:models,chap:skeleta,chap:metrics}, together with nets, weak
convergence, and extended-real arithmetic from
\Cref{chap:topological-tools}.  Relative intersection theory and the
birational input may be reviewed in
\Cref{app:birational-geometry,app:intersection-bdivisors}; multiplier ideals
and compact upper-envelope calculus are collected in
\Cref{app:ideals,app:nets}.  The local CLD theory is not a prerequisite for
the BFJ construction, but it becomes relevant when the resulting measures
are compared in \Cref{part:comparisons}.

\paragraph{Conceptual clusters.}
\label{front:part3-conceptual-clusters}
The model and compactness cluster
\Cref{chap:bfj-model-forms,chap:bfj-psh,chap:bfj-compactness} fixes model
closed forms, defines psh functions skeletally, and establishes the
compactness and approximation mechanisms used later.  Envelopes and measure
theory are organized in
\Cref{chap:bfj-envelopes,chap:bfj-ma,chap:bfj-capacity}; energy and the
variational method occupy
\Cref{chap:bfj-energy,chap:bfj-differentiability,chap:bfj-finite-energy-measures}.
Continuity and change of field are kept separate in
\Cref{chap:bfj-continuity,chap:bfj-functoriality}, and
\Cref{chap:bfj-examples} closes the part with curves, toric varieties, and a
full synthesis of the dependencies.

\paragraph{Dependency warning.}
\label{front:part3-dependency-warning}
The analytic chain rests on algebraic input: relative intersection theory,
the Hodge index theorem on models, multiplier ideals, and approximate Zariski
decompositions.  Each use records the model, characteristic, and positivity
hypotheses.  Decreasing regularization is a theorem, not part of the
definition of a BFJ psh function, and continuity of envelopes is not a formal
consequence of compactness.  Readers should also distinguish a statement for
model functions, its extension to continuous semipositive metrics, and its
finite-energy extension.

\paragraph{Reading routes.}
\label{front:part3-reading-routes}
The \emph{short route} is
\Cref{chap:bfj-model-forms,chap:bfj-psh,chap:bfj-ma,chap:bfj-energy,chap:bfj-examples};
it gives the main objects and their model tests.  The \emph{standard route}
inserts
\Cref{chap:bfj-compactness,chap:bfj-envelopes,chap:bfj-capacity} and continues
through \Cref{chap:bfj-finite-energy-measures}, thereby reaching the
variational Monge--Amp\`ere theory.  The \emph{reference route} reads the part
in order, including continuity and functoriality, with
\Cref{app:nets,app:intersection-bdivisors,app:normalizations} at hand.  A
reader chiefly interested in CLD comparison may study the model locus in
\Cref{chap:bfj-model-forms,chap:bfj-ma} and then turn to
\Cref{chap:compare-ma}.

\paragraph{Main outcomes and interfaces.}
\label{front:part3-main-outcomes-and-interfaces}
At the end of the part, the reader should know how BFJ psh functions are
assembled from models and skeleta, how mixed model intersections extend to a
Monge--Amp\`ere operator on the stated domains, how capacity controls
quasicontinuity, and how energy and dual energy organize the variational
equation.  The common model measures meet CLD currents in
\Cref{chap:compare-ma}.  The trivially valued theory in \Cref{part:bj} uses
many analogous convex and intersection-theoretic mechanisms, but replaces a
fixed ground valuation by a homogeneous valuation cone; it is not obtained by
formally setting the norm equal to one.

\paragraph{Normalization and scope.}
\label{front:part3-normalization-and-scope}
Unless explicitly normalized, the Monge--Amp\`ere measure on the bounded
and full-mass domains of an ample line bundle has total mass equal to its
top self-intersection. Energy is measured
relative to a fixed reference metric and shifts affinely under addition of a
constant.  Divisorial points are normalized by the ground valuation, and
changing that normalization rescales skeletal lengths and function values but
not algebraic intersection numbers.  Field, residue-characteristic,
smoothness, and ampleness assumptions must therefore be carried through every
comparison; \Cref{app:normalizations} records the corresponding sign and mass
tests.
